% Agent 19: Pages 415--418 (PDF pp.~36--37)
% Coverage: Section 5.2 (continued) -- angular momentum conservation,
%           mass conservation, energy conservation, heat generation,
%           surface temperature, spectrum
% Equations (5.2.12)--(5.2.23)

%%% ====================================================================
%%% SECTION 5.2  Accretion in Binary Systems: The General Picture (cont.)
%%% ====================================================================

%% --- continuation from previous agent (page 413, right column) ---

Viscous stresses in the orbiting disk gradually remove angular momentum
from each gas element, permitting it to gradually spiral inward toward the
compact star. The angular momentum removed is transported by the viscous
stresses from the inner parts of the disk to the outer parts, and is then picked
up and carried away by passing gas (see above). At the same time, the shearing
orbital motion of the gas in the inner parts of the disk, acting against the
viscous stresses, produces heat (``frictional heating''). Much more heat is
generated than the gas can store, so most of the heat gets radiated away from
the top and bottom faces of the disk.

The total energy radiated by a unit mass of gas during its passage inward
through the disk must equal, approximately,\footnote{``Approximately'' because the compact star itself can feed energy into the disk by means
of viscous stresses, and that energy can subsequently be converted to heat and radiated
away; see below.} the gravitational binding energy
of the unit mass when it reaches the inner edge of the disk, $\tilde{E}_{\text{bind}}$. For a white
dwarf or neutron star the inner edge of the disk is near the star's surface, so
\begin{equation}
\tilde{E}_{\text{bind}} \simeq \tfrac{1}{4}(M_c/R_c)
\tag{5.2.9a}
\end{equation}
\begin{equation}
\simeq 10^{-4} \quad \text{for white dwarf}.
\notag
\end{equation}
\begin{equation}
\simeq 0.05 \quad \text{for neutron star}.
\notag
\end{equation}

For a black hole the inner edge of the disk is at the last stable circular orbit,
which has
\begin{equation}
\tilde{E}_{\text{bind}} \simeq 0.057 \quad \text{for nonrotating hole}.
\tag{5.2.9b}
\end{equation}
\begin{equation}
\tilde{E}_{\text{bind}} \simeq 0.42 \quad \text{for maximally rotating hole}.
\notag
\end{equation}

Thus, in order of magnitude the total luminosity of the disk must be
\begin{equation}
L \sim 10^{-4}\,\dot{M}_0 \sim (10^{34} \text{ ergs/sec})
\left(\frac{\dot{M}_0}{10^{-9} M_\odot/\text{yr}}\right)
\quad \text{for white dwarf},
\tag{5.2.10}
\end{equation}
\begin{equation}
L \sim 0.1\,\dot{M}_0 \sim (10^{37} \text{ ergs/sec})
\left(\frac{\dot{M}_0}{10^{-9} M_\odot/\text{yr}}\right)
\quad \text{for neutron star or hole}.
\notag
\end{equation}

Here $\dot{M}_0$ is the accretion rate. In the case of a neutron star or white dwarf, the
gas hits the star (or its magnetic field) with a kinetic energy roughly equal to the
binding energy $\tilde{E}_{\text{bind}}$. In a steady state the energy liberated by this collision
must all get radiated away, so the star itself will emit a luminosity of the same
order as that emitted by the disk. In the case of a black hole, by contrast, the
disk is the only source of radiation.

If the total luminosity approaches the ``Eddington limit''
\begin{equation}
L_{\text{crit}} \simeq (1 \times 10^{38} \text{ ergs/sec})\,(M/M_\odot)
\tag{5.2.11a}
\end{equation}
(eq.~4.5.4), then radiation pressure will destroy the disk, and the general nature
of the accretion will be quite different from that depicted above. For details
see \S5.13, on the ``subcritical case.''

$L \leq L_{\text{crit}}$, with a well-defined, thin disk. Note that this subcritical case
corresponds to an accretion rate of
\begin{equation}
\dot{M}_0 \ll \dot{M}_{\text{crit}} \sim (10^{-5} M_\odot/\text{yr})(M/M_\odot)
\quad \text{for neutron star or black hole}.
\tag{5.2.11b}
\end{equation}
\begin{equation}
\sim (10^{-5} M_\odot/\text{yr})(M/M_\odot)
\quad \text{for white dwarf}
\notag
\end{equation}

For a semiquantitative, Newtonian analysis of the subcritical disk structure,
introduce an inertial frame centered on the compact star and ignore the tidal
gravitational forces of the normal star. Let the rate at which gas is deposited
into the disk, $\dot{M}_0$ (``accretion rate''), be given.

The specific angular momentum of the gas at radius $r$ is
\begin{equation}
\tilde{L} = (M/r^3)^{1/2}\,r^2 = (Mr)^{1/2},
\tag{5.2.12}
\end{equation}
where $M$---previously denoted $M_c$---is the mass of the compact star. Since the
radius of the compact star is far less than that of the outer edge of the disk, a gas
element must lose nearly all of its initial angular momentum by the time it reaches
the star. Thus, the rate at which angular momentum is removed from the disk by
passing gas must be
\begin{equation}
\dot{j} = \dot{M}_0 \times (\tilde{L} \text{ evaluated at outer edge of disk}, r_0)
= \dot{M}_0 (Mr_0)^{1/2}.
\tag{5.2.13}
\end{equation}

In a steady state the accretion rate and the outer edge of the disk will adjust
themselves until this relation is satisfied, with angular momentum continually
being removed from the outer edge by passing gas.

Return to the inner regions of the disk. Let $2h$
be the disk thickness, and let $\Sigma = 2h\rho_0$ be the surface density (g/cm$^2$),
$\rho_0$ be the mass density (g/cm$^3$) at radius
$r$. Let $v^{\hat{r}}$ be the radial velocity of the gas. (Note: $v^{\hat{r}} < 0$.) The orbital velocity
and angular velocity $\Omega$ are
\begin{equation}
v^{\hat{\phi}} = \Omega r = (M/r)^{1/2}\,.
\tag{5.2.14}
\end{equation}

The viscosity will never get strong enough to allow large radial velocities; i.e.,
$|v^{\hat{r}}|$ will always be far smaller than $v^{\hat{\phi}}$. Let $t_{\hat{r}\hat{\phi}}$ be the viscous stress (components
relative to an orthonormal frame at radius $r$). This stress is related to the shear
of the circular Keplerian orbits by
\begin{equation}
t_{\hat{r}\hat{\phi}} = -2\eta\,\sigma_{\hat{r}\hat{\phi}}
= -\tfrac{3}{2}\eta\,(M/r^3)^{1/2},
\tag{5.2.16}
\end{equation}
where $\eta$ is the coefficient of dynamic viscosity (see \S2.5). The sources of
viscosity will be discussed below. Let $F$ be the flux of radiation (ergs/cm$^2$ sec)

off the upper face of the disk. An equal flux $F$ will come off the lower face.

The steady-state structure of the disk is governed by four conservation laws
(conservation of mass, angular momentum, energy, and vertical momentum);
by the nature of the viscosity---i.e., value of $\eta$---and by the law of radiative transfer
from the inside of the disk to its surface.

\textit{Rest-mass conservation}

Mass must flow across a cylinder of radius $r$ at a rate equal to the accretion rate
\begin{equation}
-2\pi r \cdot 2h \cdot \rho_0\, v^{\hat{r}} = \dot{M}_0\,.
\tag{5.2.15}
\end{equation}

%%% --- Angular-momentum conservation ---

\textit{Angular-momentum conservation}

The rate at which angular momentum is carried inward across radius $r$ by
inflowing gas must equal the rate at which viscous stresses carry angular
momentum out of radius $r$, plus the rate $\dot{j}_c$ at which angular momentum is
deposited in the compact star
\begin{equation}
\dot{M}_0 (Mr)^{1/2} = 2\pi r \cdot 2h \cdot t_{\hat{r}\hat{\phi}} \cdot r + \dot{j}_c\,.
\tag{5.2.17}
\end{equation}

The specific angular momentum $\tilde{L}_c$ deposited into the compact star cannot
exceed the Keplerian angular momentum at the inner edge of the disk $r_I$;
hence
\begin{equation}
\dot{j}_c = \beta \dot{M}_0 (Mr_I)^{1/2} \quad \text{for some} \quad |\beta| \leq 1.
\notag
\end{equation}

Solving for the product of stress and disk thickness, we obtain
\begin{equation}
2h t_{\hat{r}\hat{\phi}} = \frac{\dot{M}_0}{2\pi r^2}\left[(Mr)^{1/2}
- \beta(Mr_I)^{1/2}\right].
\tag{5.2.18}
\end{equation}
\begin{equation}
\simeq \frac{\dot{M}_0(Mr)^{1/2}}{2\pi r^2}
\quad \text{for} \quad r \gg r_I.
\notag
\end{equation}

Notice that in a steady state the product $2h t_{\hat{r}\hat{\phi}}$ is determined uniquely by the
accretion rate and the mass of the hole. If, at some radius, $2h t_{\hat{r}\hat{\phi}}$ is (temporarily)
smaller than the steady-state value, then the viscous stresses there will not be sufficient to
handle the mass flow $\dot{M}_0$. As a result, mass will accumulate in the deviant region,
increasing thereby the disk thickness $h$ and/or the stresses $t_{\hat{r}\hat{\phi}}$ there, and returning
the deviant region to a steady-state structure. For a quantitative analysis of this
``approach to steady state'' see Thorne~(1973b)\cite{Thorne1973b}.

%%% --- Energy conservation ---

\textit{Energy conservation}

Heat is generated in the disk by viscosity at a rate, per unit volume, given by
\begin{equation}
\varepsilon = (2h t_{\hat{r}\hat{\phi}})(-2\sigma_{\hat{r}\hat{\phi}})
\notag
\end{equation}
(see \S2.5). Thus, the heat generated per unit area is
\begin{equation}
2h\varepsilon = (2h t_{\hat{r}\hat{\phi}})(-2\sigma_{\hat{r}\hat{\phi}})
= \frac{3\dot{M}_0}{4\pi r^2}\,\frac{M}{r}
\left[1 - \beta\!\left(\frac{r_I}{r}\right)^{1/2}\right];
\tag{5.2.19}
\end{equation}
and the heat generated between radii $r_1$ and $r_2$ is
\begin{equation}
\int_{r_1}^{r_2} 2h\varepsilon\cdot 2\pi r\,dr
= \tfrac{3}{2}\dot{M}_0
\left[\frac{M}{r_1}
- \frac{2\beta}{3}\left(\frac{r_I}{r_1}\right)^{1/2}\!\frac{M}{r_1}
- \frac{M}{r_2}
+ \frac{2\beta}{3}\left(\frac{r_I}{r_2}\right)^{1/2}\!\frac{M}{r_2}
\right]
\notag
\end{equation}
\begin{equation}
\simeq \tfrac{3}{2}\dot{M}_0
\left(\frac{M}{r_1} - \frac{M}{r_2}\right)
\quad \text{for} \quad r_2 > r_1 \gg r_I.
\tag{5.2.20}
\end{equation}

It is instructive to examine the origin of this heat energy by asking at what
net rate energy is being deposited in the region $r_1 < r < r_2$. Gravitational
potential energy is released at a rate $\dot{M}_0(M/r_1 - M/r_2)$, but only half of this
energy can go into orbital kinetic energy
(virial theorem). Thus, the rate at which gravitational energy gets converted
into heat is
\begin{equation}
\tfrac{1}{2}\dot{M}_0(M/r_1 - M/r_2).
\notag
\end{equation}

The viscous stresses transport outward not only angular momentum, but also
energy. The rate at which energy is transported across radius $r$ is
\begin{equation}
\dot{E} = \Omega J = \dot{M}_0(M/r)\left[1 - \beta(r_I/r)^{1/2}\right].
\notag
\end{equation}

Thus, energy is deposited by viscosity between radii $r_1$ and $r_2$ at a rate
\begin{equation}
\dot{M}_0
\left[\frac{M}{r_1}\left(1 - \beta\!\left(\frac{r_I}{r_1}\right)^{1/2}\right)
- \frac{M}{r_2}\left(1 - \beta\!\left(\frac{r_I}{r_2}\right)^{1/2}\right)
\right].
\notag
\end{equation}

This energy deposition rate plus the deposition rate for gravitational energy is
equal to the total heating rate. Notice that at radii $r \gg r_I$, gravity accounts
directly for only one-third of the heat. The remaining two-thirds is transported
into the heating region by viscous stresses. Much of the early literature on disk
accretion, e.g., Lynden-Bell~(1969)\cite{LyndenBell1969}, failed to take account of energy transport
by viscous stresses, and therefore underestimated by a factor of 3 the heating at
$r \gg r_I$.

If none of the heat were radiated away, the thermal energy of the disk
would be 3/2 times its gravitational potential energy; and its temperature
would thus be
\begin{equation}
T \sim (M/r) m_p \sim (10^{13}\,\text{K})(r/10^5\,\text{cm})^{-1}(M/M_\odot).
\notag
\end{equation}

Such temperatures are absurdly high. Thermal bremsstrahlung and other
radiative processes will never permit such temperatures to arise. Instead, they
will quickly remove almost all of the heat that is generated. This means that

the total flux $2F$ from the top and bottom faces of the disk must equal the
heating rate per unit area, as calculated above (the heat flows out of the disk
vertically rather than radially because the disk is so thin):
\begin{equation}
F = \frac{3\dot{M}_0}{8\pi r^2}\,\frac{M}{r}
\left[1 - \beta\!\left(\frac{r_I}{r}\right)^{1/2}\right].
\tag{5.2.21}
\end{equation}

The total power radiated is thus
\begin{equation}
L = \int_{r_I}^{\infty} 2F \cdot 2\pi r\,dr
= (\tfrac{3}{2} - \beta)\dot{M}_0\,\frac{M}{r_I}\,.
\tag{5.2.22}
\end{equation}

Of this total power radiated, $\tfrac{1}{2}\dot{M}_0 M/r_I$ comes from gravity, while $(1 - \beta)\dot{M}_0 M/r_I$
comes from the rotational energy of the star.

\textit{Spectrum}

A detailed treatment of the spectrum radiated will be given in \S5.10. Here we
only note that, if the radiation is blackbody, then the surface temperature of
the disk at radius $r$ must be
\begin{equation}
T_s = \left(\frac{4F}{\sigma}\right)^{1/4}
\simeq (3 \times 10^7\,\text{K})
\left(\frac{\dot{M}_0}{10^{-9}\,M_\odot/\text{yr}}\right)^{1/4}
\!\left(\frac{M}{M_\odot}\right)^{-1/2}
\!\left(\frac{r}{r_I}\right)^{-3/4}
\left[1 - \beta\!\left(\frac{r_I}{r}\right)^{1/2}\right]^{1/4}\!.
\tag{5.2.23}
\end{equation}

Because most of the radiation is emitted from $r/M \sim 10$ for a black hole or
neutron star, and from $r/M \sim 10^4$ for a dwarf, and because the black-body
spectrum peaks at a photon energy of
\begin{equation}
h\nu_{\max} \simeq (2.44 \times 10^{-4}\;\text{eV})(T_s/{}^{\circ}\text{K}),
\notag
\end{equation}
the spectrum of the total radiation from the disk should peak at
\begin{equation}
h\nu_{\max} \simeq (1\;\text{keV})
\left(\frac{\dot{M}_0}{10^{-9} M_\odot/\text{yr}}\right)^{1/4}
\left(\frac{M}{M_\odot}\right)^{-1/2}
\quad \text{for neutron star or hole}
\notag
\end{equation}
\begin{equation}
h\nu_{\max} \simeq (0.01\;\text{keV})
\quad \text{for white dwarf.}
\notag
\end{equation}

Thus, for reasonable accretion rates the disk can emit strongly in the X-ray
region ($\sim$1 to 10~keV); but the X-ray spectrum should fall fairly rapidly with
increasing energy. It actually turns out that electron-scattering opacity impedes
the emission of blackbody radiation, and causes the spectrum to peak at energies
higher than that (5.2.23). (See \S5.10.) However, the above estimates are still
accurate to within a factor $\sim$10.

%%% --- Vertical pressure balance ---

\textit{Vertical pressure balance}

The thickness of the disk is governed by a balance between vertical pressure force
and the tidal gravitational force of the compact star (``vertical momentum
conservation''):
\begin{equation}
\frac{dp}{dz} = \rho_0 \times \left(\text{``acceleration of gravity''}\right) = \rho_0\,\frac{Mz}{r^3}\,.
\tag{5.2.24}
\end{equation}

The approximate solution to this equation is
\begin{equation}
h \simeq (p/\rho_0)^{1/2}(r^3/M)^{1/2} = c_s/\Omega.
\tag{5.2.25}
\end{equation}

Here $h$ is the half-thickness of the disk, $c_s$ is the speed of sound in the gas, and
$\Omega = (M/r^3)^{1/2}$ is the angular velocity of the gas.
